% ===================================================================
%  AntennaForge — Field Guide
%  Parameter Effects and Implementation Details
% ===================================================================
\documentclass[11pt,a4paper]{article}

% ── Packages ──────────────────────────────────────────────────────
\usepackage[utf8]{inputenc}
\usepackage[T1]{fontenc}
\usepackage{lmodern}
\usepackage[margin=1in]{geometry}
\usepackage{xcolor}
\usepackage{hyperref}
\usepackage{booktabs}
\usepackage{longtable}
\usepackage{amsmath,amssymb}
\usepackage{titlesec}
\usepackage{fancyhdr}
\usepackage{parskip}
\DeclareMathOperator{\sinc}{sinc}

% ── Colour Definitions ───────────────────────────────────────────
\definecolor{afblue}{HTML}{3B82F6}
\definecolor{afdark}{HTML}{0F172A}
\definecolor{afgray}{HTML}{475569}
\definecolor{afred}{HTML}{DC2626}

% ── Hyperref Setup ───────────────────────────────────────────────
\hypersetup{
    colorlinks=true,
    linkcolor=afblue,
    urlcolor=afblue,
    pdftitle={AntennaForge Field Guide},
    pdfauthor={AntennaForge Project},
}

% ── Header / Footer ─────────────────────────────────────────────
\pagestyle{fancy}
\fancyhf{}
\fancyhead[L]{\bfseries AntennaForge Field Guide}
\setlength{\headheight}{13.6pt}
\fancyhead[R]{\thepage}
\renewcommand{\headrulewidth}{0.4pt}

% ── Section Formatting ───────────────────────────────────────────
\titleformat{\section}{\Large\bfseries\color{afdark}}{}{0em}{\thesection\quad}
\titleformat{\subsection}{\large\bfseries\color{afblue}}{}{0em}{\thesubsection\quad}

\begin{document}

\begin{center}
    {\Huge\bfseries\color{afdark} AntennaForge Field Guide}\\[0.5em]
    {\Large Parameter Effects \& Implementation Logic}\\[1em]
    {\large\today}
\end{center}

\vspace{1cm}

\tableofcontents
\newpage

\section{Introduction}

This Field Guide complements the User Guide by providing a detailed technical explanation of how GUI parameters influence the generated radiation patterns. It focuses on the mathematical models and code logic underlying each control.

\begin{quote}
    \textbf{\color{afred} DISCLAIMER:} AntennaForge generates \emph{synthetic} radiation patterns based on analytical mathematical models (e.g., Gaussian, $\cos^n$, sinc, Jinc). These are idealised approximations designed for system-level simulation. They are \textbf{not} derived from electromagnetic simulation of physical geometry (Method of Moments, FEM, FDTD) and may not capture near-field effects, complex structural scattering, or specific hardware imperfections unless explicitly modelled via the available physics features.
\end{quote}

\newpage
\section{Core Parameters}

\subsection{Antenna Type}
The \textbf{Antenna Type} selection determines the fundamental diffraction kernel used to shape the main beam. Changing this switches the mathematical function $F(\theta)$ called by the engine.

\textbf{Magnitude of Effect:} \textbf{Drastic}. This is the primary determinant of pattern shape.

\begin{longtable}{p{3cm} p{12cm}}
\toprule
\textbf{Type} & \textbf{Implementation Logic \& Effect} \\
\midrule
\textbf{LPDA} & Uses an \textbf{Elliptical Gaussian} model ($G \propto 0.5^{r^2}$). This produces a smooth, bell-shaped main beam typical of log-periodic arrays. It lacks the distinct nulls of aperture antennas unless sidelobes are explicitly enabled. \\
\textbf{Dish} & Uses the \textbf{Jinc squared} (Airy) function ($G \propto [2J_1(u)/u]^2$). This models a circular aperture. It produces a narrow main beam with a first sidelobe at $-17.6$\,dB (theoretical). \\
\textbf{Panel} & Uses \textbf{Sinc squared} functions ($G \propto \sinc ^2(u_{az}) \cdot \sinc^2(u_{el})$). Models a rectangular aperture with uniform illumination. Produces first sidelobes at $-13.3$\,dB. \\
\textbf{Horn} & Hybrid model: \textbf{Sinc squared} in E-plane (uniform distribution) and \textbf{Cosine-tapered} in H-plane. This captures the asymmetry of pyramidal horns. \\
\textbf{Omni} & \textbf{Cosine power} ($G \propto \cos^n(\theta_{el})$) in elevation; uniform in azimuth. $n$ is derived from the elevation beamwidth. \\
\textbf{Monopole} & \textbf{Dipole far-field} formula ($G \propto [\cos(kL/2 \cdot \sin\theta_{el}) - \cos(kL/2)]^2 / [\cos(\theta_{el})]^2$). Omnidirectional in azimuth (toroidal). Element length is configurable in wavelengths ($0.25\lambda$ quarter-wave, $0.5\lambda$ half-wave). \\
\textbf{Array} & \textbf{Phasor Summation}. Computes $\sum w_n e^{j(k d \sin \theta + \beta)}$ with selectable element patterns (Gaussian, cosine, or patch). Supports linear, planar, and circular geometries with optional mutual coupling and per-element amplitude/phase errors. This is the only model that naturally produces grating lobes if element spacing $d > 0.5\lambda$. \\
\bottomrule
\end{longtable}

\subsection{Max Gain (dBi)}
This is a scalar multiplier applied at the very end of the pattern computation.
\begin{itemize}
    \item \textbf{Magnitude:} \textbf{Linear Scaling}. Shifts the entire curve up/down.
    \item \textbf{Code Implementation:} The engine generates a normalized pattern shape $F(\theta)$ (where $\max(F) \approx 1.0$) and then scales it:
    \[ G(\theta) = F(\theta) \cdot 10^{(\text{MaxGain}/10)} \]
    \item \textbf{Effect:} Sets the absolute level of the pattern peak. It acts as a \textbf{scalar multiplier}, not a hard clipping limit (the pattern shape is preserved).
    \item \textbf{Independence:} Max Gain is \textbf{decoupled} from beamwidth and FTB (unless Gain-Beamwidth Coupling is enabled). This allows you to specify the exact peak gain from a datasheet without the tool automatically adjusting it to satisfy the conservation of energy integral ($4\pi$ steradians).
\end{itemize}

\subsection{Front-to-Back Ratio (dB)}
Controls the floor of the rear hemisphere.
\begin{itemize}
    \item \textbf{Magnitude:} \textbf{Noticeable} (Rear Hemisphere).
    \item \textbf{Code:} The engine computes a \texttt{back\_lobe\_level} based on this ratio.
    \item \textbf{Logic:} The pattern is a blend of the front model and the back model.
    \item \textbf{Effect:} Increasing this value pushes the pattern level at $\pm 180^\circ$ downwards. It does not affect the main beam shape.
    \item \textbf{Physics Note:} In a physical antenna, increasing the backlobe (lowering FTB) would reduce the forward gain due to conservation of energy. In AntennaForge, these are independent knobs. You can increase the backlobe level without penalising the main beam gain, allowing for "worst-case" interference modelling where both forward and backward radiation are high.
\end{itemize}

\newpage
\section{Beamwidth \& Frequency Scaling}

\subsection{Azimuth / Elevation Beamwidth}
These inputs define the $-3$\,dB points of the pattern at the \textbf{Reference Frequency}.
\begin{itemize}
    \item \textbf{Magnitude:} \textbf{Drastic}. Defines the width of the main beam.
    \item \textbf{Code:} Converted into a scaling factor $k$ or exponent $n$.
    \item \textbf{Effect:}
    \begin{itemize}
        \item For $\cos^n$ models (Omni, LPDA backlobe): Defines how "pointy" the cosine function is.
        \item For Aperture models (Dish, Panel, Horn): Defines the coordinate scaling factor $u$.
    \end{itemize}
\end{itemize}

\subsection{Frequency-Dependent Beamwidth}
Real antennas have fixed physical apertures, so their electrical size (in wavelengths) increases with frequency, causing the beam to narrow.

\begin{itemize}
    \item \textbf{Exponent Mode:} $BW(f) = BW_{ref} \cdot (f_{ref} / f)^\alpha$.
    \item \textbf{Percentage Mode:} $BW(f) = BW_{ref} \cdot (1 - \frac{pct}{100})^{\log_2(f/f_{ref})}$.
    \item \textbf{Three-Point Mode:} Fits a quadratic curve through user-defined BWs at $f_{min}, f_{mid}, f_{max}$.
\end{itemize}

\textbf{Magnitude:} \textbf{Noticeable}.
\textbf{Effect on Pattern:}
If enabled, you will see the main beam get narrower (sharper) in higher frequency slices. If disabled, the beam shape remains identical across all files.

\newpage
\section{Physics Features}

\subsection{Sidelobes (Taylor)}
Replaces the theoretical nulls/lobes of the kernel with a controlled envelope.
\begin{itemize}
    \item \textbf{Magnitude:} \textbf{Noticeable}. Changes the pattern from a simple bell curve to a structured diffraction pattern.
    \item \textbf{Parameters:} First SLL (dB), Decay Rate (dB/lobe), Count.
    \item \textbf{Code Implementation:} The code calculates a normalized angle $u$. If $u$ is outside the main beam, it switches to a sidelobe function:
    \[ G_{sl} = A_{1st} \cdot \text{decay}^{lobe\_index} \cdot \cos^2(\text{phase}) \]
    \item \textbf{Effect:}
    \begin{itemize}
        \item \textbf{First SLL:} Sets the height of the "shoulders" next to the main beam.
        \item \textbf{Decay:} Controls how fast the ripples fade away. High decay = cleaner pattern.
    \end{itemize}
\end{itemize}

\subsection{Pattern Breakup}
Simulates the "messiness" of real measurements at high off-axis angles.
\begin{itemize}
    \item \textbf{Magnitude:} \textbf{Subtle to Moderate}. Adds texture to the noise floor.
    \item \textbf{Code:} Adds a pseudo-random ripple function:
    \[ \Delta_{dB} = A \cdot \sin(D \cdot \theta) \cdot \text{random\_factor} \]
    \item \textbf{Effect:} Adds "grass" or "fuzz" to the pattern floor, typically starting at $\pm 90^\circ$. Useful for testing robust peak-finding algorithms.
\end{itemize}

\subsection{Ground Reflection}
Models a 2-ray multipath environment (Direct + Reflected).
\begin{itemize}
    \item \textbf{Magnitude:} \textbf{Drastic} (Elevation only). Creates deep interference nulls.
    \item \textbf{Code:} Computes phase difference $\Delta \phi = 4\pi h \sin(\theta_{el}) / \lambda$.
    \item \textbf{Factor:} $F = \sqrt{1 + |\Gamma|^2 + 2|\Gamma|\cos(\Delta \phi + \pi)}$.
    \item \textbf{Effect:} Creates deep nulls and peaks in the \textbf{Elevation} pattern. The number of lobes increases with antenna height $h$.
\end{itemize}

\subsection{Cross-Polarisation}
Generates a secondary pattern file representing energy leaked into the orthogonal plane.
\begin{itemize}
    \item \textbf{Magnitude:} \textbf{Secondary Output}. Generates a separate file; usually 15-30dB down.
    \item \textbf{Code:} Modelled as a "Maltese Cross" shape (peaks at $\pm 45^\circ$ azimuth/elevation diagonals) plus a boresight leakage term.
    \item \textbf{Effect:} The \texttt{\_xpol.csv} file will show a pattern with a null at boresight (if isolation is good) and peaks in the diagonal quadrants.
\end{itemize}

\subsection{VSWR Rolloff}
Simulates impedance mismatch at the band edges.
\begin{itemize}
    \item \textbf{Magnitude:} \textbf{Subtle}. Only affects gain at the very edges of the frequency band.
    \item \textbf{Code:} Calculates a loss factor based on how close $f$ is to $f_{min}$ or $f_{max}$.
    \item \textbf{Effect:} The \textbf{Gain vs Frequency} plot will show a flat top with rounded "shoulders" dropping off at the start and end of the band.
\end{itemize}

\subsection{Asymmetry / Squint}
\begin{itemize}
    \item \textbf{Magnitude:} \textbf{Subtle}. Small angular shifts.
    \item \textbf{Code:} Adds an offset to the angle query: $G(\theta - \theta_{squint})$.
    \item \textbf{Effect:} Shifts the main beam peak away from $0^\circ$. Useful for testing pointing error budgets.
\end{itemize}

\newpage
\section{Simulation Settings}

\subsection{Sigmoid K (Front-Fade Steepness)}
Controls the blending region between the front hemisphere (main beam) and the rear hemisphere (back lobe).
\begin{itemize}
    \item \textbf{Magnitude:} \textbf{Subtle}. Affects the blend region at $\pm 90^\circ$.
    \item \textbf{Code:} Weight $w = 1 / (1 + e^{-k \cos \theta})$.
    \item \textbf{Effect:}
    \begin{itemize}
        \item \textbf{High K (e.g., 40):} Sharp transition at $\pm 90^\circ$. Can look artificial (like a hard baffle).
        \item \textbf{Low K (e.g., 12):} Smooth transition. More realistic for wire antennas or small apertures.
    \end{itemize}
\end{itemize}

\subsection{Noise}
\begin{itemize}
    \item \textbf{Magnitude:} \textbf{Subtle}. Adds jitter to the curve.
    \item \textbf{Uniform:} Adds random value in $\pm$ Range.
    \item \textbf{Gaussian:} Adds random value with $\sigma =$ Value.
    \item \textbf{Quantization:} Rounds dB values to nearest Step.
    \item \textbf{Effect:} Makes the pattern curves "jittery". Quantization creates a stair-step effect visible in cuts.
\end{itemize}

\newpage
\section{Type-Specific Parameters}

In addition to the core parameters shared by all types (gain, beamwidth, FTB, sigmoid~K), several antenna types have unique parameters that control type-specific physics.

\subsection{Panel: Mechanical Tilt}
\begin{itemize}
    \item \textbf{Config key:} \texttt{mechanical\_tilt\_deg} (replaces legacy \texttt{panel\_downtilt\_deg})
    \item \textbf{Default:} 0.0
    \item \textbf{Range:} $-45^\circ$ to $+45^\circ$ (positive = downtilt)
    \item \textbf{Effect:} Offsets the elevation angle before pattern computation: $\theta'_{el} = \theta_{el} - \text{tilt}$. Also available for Horn, Dish, and Array types.
\end{itemize}

\subsection{Dish: Feed Edge Taper}
\begin{itemize}
    \item \textbf{Config key:} \texttt{feed\_edge\_taper\_db}
    \item \textbf{Default:} $-12.0$\,dB
    \item \textbf{Range:} $-30$\,dB to $0$\,dB
    \item \textbf{Effect:} Applies a Gaussian taper envelope $\exp(-r^2 / 2\sigma^2)$ to the Airy pattern. More negative values suppress sidelobes but widen the main beam. At $0$\,dB (uniform illumination), the Airy pattern sidelobes are at their theoretical maximum ($-17.6$\,dB first SLL).
    \item \textbf{Physics:} Models the illumination taper from the feed horn at the dish edge. Real dishes typically have $-10$ to $-15$\,dB edge taper.
\end{itemize}

\subsection{Horn: Aperture Wavelengths}
\begin{itemize}
    \item \textbf{Config key:} \texttt{horn\_aperture\_wavelengths}
    \item \textbf{Default:} \texttt{None} (manual beamwidth mode)
    \item \textbf{Effect:} When set, the beamwidth is derived from the physical aperture size and scales automatically with frequency:
    \[ BW_H(f) = \frac{67^\circ}{\text{aperture} \cdot f/f_{ref}}, \qquad BW_E(f) = \frac{51^\circ}{\text{aperture} \cdot f/f_{ref}} \]
    This models the physical relationship between aperture size and beamwidth. When \texttt{None}, beamwidths are set manually via the standard beamwidth controls.
\end{itemize}

\subsection{Monopole: Physical Length}
\begin{itemize}
    \item \textbf{Config key:} \texttt{monopole\_physical\_length\_m}
    \item \textbf{Default:} \texttt{None} (fixed electrical length mode)
    \item \textbf{Effect:} When set, the electrical length scales with frequency:
    \[ L(f) = \frac{\text{physical\_length}}{c / f} \]
    This models a fixed physical element whose pattern evolves across the band (e.g., a $0.25\lambda$ monopole at $f_{min}$ becomes $0.5\lambda$ at $2 \times f_{min}$, causing the elevation pattern to narrow and eventually develop nulls). When \texttt{None}, every frequency slice uses the same \texttt{element\_length\_wavelengths} value.
\end{itemize}

\subsection{Array: Element Pattern}
\begin{itemize}
    \item \textbf{Config key:} \texttt{array\_element\_pattern}
    \item \textbf{Default:} \texttt{gaussian}
    \item \textbf{Options:}
    \begin{itemize}
        \item \texttt{gaussian} --- $G_{elem} = 0.5^{r^2}$ where $r^2 = (\theta_{az}/\theta_{hw})^2 + (\theta_{el}/\theta_{hw})^2$. Smooth, no nulls.
        \item \texttt{cosine} --- $G_{elem} = [\cos(\pi/2 \cdot \theta/\theta_{3dB})]^2$ per plane. Physically motivated for dipole-like elements.
        \item \texttt{patch} --- $G_{elem} = \cos^{1.5}(\theta_{az}) \cdot \cos^{1.5}(\theta_{el})$. Empirical model for microstrip patch elements.
    \end{itemize}
    \item \textbf{Effect:} Changes the envelope of the total array pattern. Cosine and patch elements roll off faster than Gaussian, producing a narrower visible-region pattern.
\end{itemize}

\subsection{Array: Element Errors}
\begin{itemize}
    \item \textbf{Config keys:} \texttt{array\_rms\_amplitude\_error\_db}, \texttt{array\_rms\_phase\_error\_deg}
    \item \textbf{Defaults:} 0.0 (ideal)
    \item \textbf{Effect:} Applies Gaussian-distributed random perturbations to each element's amplitude weight and phase offset. Uses a deterministic seed for repeatability. Models manufacturing tolerances, mounting errors, and feed network imperfections.
    \item \textbf{Impact:} Raises sidelobe levels, fills in nulls, and slightly reduces peak gain. Typical values: $0.3$--$1.0$\,dB amplitude, $3$--$10^\circ$ phase.
\end{itemize}

\subsection{Array: Circular Orientation}
\begin{itemize}
    \item \textbf{Config key:} \texttt{array\_circular\_orientation}
    \item \textbf{Default:} \texttt{parallel}
    \item \textbf{Options:}
    \begin{itemize}
        \item \texttt{parallel} --- All elements point in the same direction (broadside). The array factor provides directivity; element patterns are identical.
        \item \texttt{radial} --- Each element points radially outward from the array centre. Gain at angle $\phi$ is weighted by $\cos(\phi - \phi_n)$ for element $n$. Produces more uniform azimuth coverage.
    \end{itemize}
\end{itemize}

\newpage
\section{Common Adjustment Scenarios}

This section provides a quick reference for achieving specific pattern shaping goals.

\subsection{I want to change the Peak Gain...}

\textbf{...without affecting the beamwidth (Parametric Mode)}
\begin{itemize}
    \item \textbf{Action:} Change the \textbf{Max Gain (dBi)} field.
    \item \textbf{Mechanism:} AntennaForge applies gain as a scalar multiplier after the shape is generated.
    \item \textbf{Use Case:} Matching a specific datasheet value where you want exact control over both gain and beamwidth independently.
\end{itemize}

\textbf{...and have the beamwidth change naturally (Physics Mode)}
\begin{itemize}
    \item \textbf{Action:} Enable \textbf{Gain-Beamwidth Coupling} in the Physics Features card. Select "Gain drives beamwidth".
    \item \textbf{Mechanism:} Changing the Max Gain will automatically recalculate the beamwidths based on the approximate directivity formula $G \approx 30,000 / (\theta_{az} \cdot \theta_{el})$.
    \item \textbf{Use Case:} Exploring trade-offs for a fixed aperture efficiency.
\end{itemize}

\subsection{I want to change the Beamwidth...}

\textbf{...without affecting the peak gain}
\begin{itemize}
    \item \textbf{Action:} Adjust the \textbf{Azimuth BW} or \textbf{Elevation BW} sliders.
    \item \textbf{Mechanism:} The shape of the main lobe is stretched or compressed, but the peak is re-normalized to \textbf{Max Gain}.
    \item \textbf{Use Case:} Fitting a pattern to a specific coverage requirement (e.g., "I need 60-degree sector coverage at 15 dBi").
\end{itemize}

\textbf{...and have the gain change naturally}
\begin{itemize}
    \item \textbf{Action:} Enable \textbf{Gain-Beamwidth Coupling}. Select "Beamwidth drives gain".
    \item \textbf{Mechanism:} Narrowing the beamwidth will increase the calculated peak gain; widening it will reduce the gain.
\end{itemize}

\subsection{I want to change the Backlobe Level...}

\textbf{...without affecting forward gain}
\begin{itemize}
    \item \textbf{Action:} Adjust \textbf{Front-to-Back (dB)}. Ensure \textbf{FTB Affects Gain} is \textbf{OFF}.
    \item \textbf{Result:} The rear hemisphere level moves up/down independently.
\end{itemize}

\textbf{...and conserve total power (LPDA only)}
\begin{itemize}
    \item \textbf{Action:} Enable \textbf{FTB Affects Gain}.
    \item \textbf{Result:} Decreasing FTB (raising the backlobe) will reduce the forward gain to account for the energy leaking out the back.
\end{itemize}

\subsection{I want to make the beam narrower at high frequencies...}
\begin{itemize}
    \item \textbf{Action:} Ensure \textbf{Freq-dependent BW (Az/El)} is \textbf{ON}.
    \item \textbf{Mechanism:} Scales beamwidth inversely with frequency ($BW \propto 1/f$).
\end{itemize}

\subsection{I want to keep the beamwidth constant across the band...}
\begin{itemize}
    \item \textbf{Action:} Turn \textbf{Freq-dependent BW (Az/El)} \textbf{OFF}.
    \item \textbf{Mechanism:} Uses the reference beamwidth for all frequency slices.
\end{itemize}

\subsection{I want to add realistic sidelobes...}
\begin{itemize}
    \item \textbf{Action:} Turn \textbf{Sidelobes (Taylor)} \textbf{ON}. Adjust \textbf{First SLL} and \textbf{Decay Rate}.
    \item \textbf{Mechanism:} Replaces the theoretical nulls with a Taylor window envelope.
\end{itemize}

\subsection{I want to remove all sidelobes (ideal main beam only)...}
\begin{itemize}
    \item \textbf{Action:} Turn \textbf{Sidelobes (Taylor)} \textbf{OFF}.
    \item \textbf{Mechanism:} Uses the pure diffraction kernel (Gaussian, sinc, cosine) which may have nulls but no added envelope structure.
\end{itemize}

\subsection{I want to simulate ground bounce...}
\begin{itemize}
    \item \textbf{Action:} Turn \textbf{Ground Reflection} \textbf{ON}. Set \textbf{Height} and \textbf{Reflection Coeff}.
    \item \textbf{Mechanism:} Adds a multipath interference term to the elevation pattern.
\end{itemize}

\subsection{I want to see cross-polarization leakage...}
\begin{itemize}
    \item \textbf{Action:} Turn \textbf{Cross-Polarization} \textbf{ON}.
    \item \textbf{Result:} Generates additional \texttt{\_xpol.csv} files.
\end{itemize}

\subsection{I want to simulate manufacturing tolerances in a phased array...}
\begin{itemize}
    \item \textbf{Action:} Select \textbf{Array} type. Set \textbf{RMS Amplitude Error} (e.g., 0.5\,dB) and/or \textbf{RMS Phase Error} (e.g., 5\textdegree).
    \item \textbf{Mechanism:} Applies Gaussian-distributed random perturbations to each element's weight and phase.
    \item \textbf{Effect:} Raises sidelobe floor, fills nulls, and slightly reduces peak gain --- matching real-world array behaviour.
\end{itemize}

\subsection{I want to simulate manufacturing alignment errors...}
\begin{itemize}
    \item \textbf{Action:} Turn \textbf{Asymmetry / Squint} \textbf{ON}. Set \textbf{Squint} and \textbf{Tilt}.
    \item \textbf{Mechanism:} Offsets the beam peak and adds random perturbations.
\end{itemize}

\subsection{I want to simulate impedance mismatch at band edges...}
\begin{itemize}
    \item \textbf{Action:} Turn \textbf{VSWR Rolloff} \textbf{ON}.
    \item \textbf{Mechanism:} Reduces gain at $f_{min}$ and $f_{max}$.
\end{itemize}

\subsection{I want to make the pattern look "messy" or measured...}
\begin{itemize}
    \item \textbf{Action:} Turn \textbf{Pattern Breakup} \textbf{ON} or increase \textbf{Noise}.
    \item \textbf{Mechanism:} Adds ripple at wide angles or random noise to the whole pattern.
\end{itemize}

\subsection{I want to sharpen the front-to-back transition...}
\begin{itemize}
    \item \textbf{Action:} Increase \textbf{Sigmoid K} (e.g., to 40 or 60).
    \item \textbf{Effect:} Creates a harder "wall" at $\pm 90^\circ$.
\end{itemize}

\subsection{I want to soften the front-to-back transition...}
\begin{itemize}
    \item \textbf{Action:} Decrease \textbf{Sigmoid K} (e.g., to 10 or 12).
    \item \textbf{Effect:} Creates a gradual blend between front and back lobes.
\end{itemize}

\subsection{I want to steer the beam (Phased Array)...}
\begin{itemize}
    \item \textbf{Action:} Select \textbf{Array} type. Adjust \textbf{Steer Az} and \textbf{Steer El}.
    \item \textbf{Mechanism:} Applies progressive phase shift to array elements.
\end{itemize}

\subsection{I want to reduce array sidelobes...}
\begin{itemize}
    \item \textbf{Action:} Set \textbf{Weighting} to \textbf{Taylor}. Adjust \textbf{Taper SLL}.
    \item \textbf{Mechanism:} Applies amplitude tapering to elements (reduces gain slightly).
\end{itemize}

\subsection{I want to see grating lobes...}
\begin{itemize}
    \item \textbf{Action:} Set \textbf{Spacing} $> 0.5\lambda$ (e.g., 0.8). Steer the beam.
    \item \textbf{Mechanism:} Spatial aliasing in the array factor.
\end{itemize}

\subsection{I want to change the frequency resolution...}
\begin{itemize}
    \item \textbf{Action:} Change \textbf{Number of Slices}.
    \item \textbf{Result:} Generates more CSV files across the band.
\end{itemize}

\subsection{I want to change the angular resolution...}
\begin{itemize}
    \item \textbf{Action:} Change \textbf{Azimuth Step} or \textbf{Elevation Step}.
    \item \textbf{Result:} Finer grid in the CSV file (larger file size).
\end{itemize}

\newpage
\section{Troubleshooting Pattern Artifacts}

\subsection{Why is there a visible line or step at $\pm 90^\circ$?}
\begin{itemize}
    \item \textbf{Cause:} The \textbf{Sigmoid K} parameter (Front-Fade Steepness) is set too high (e.g., default 40), creating a sharp transition between the front and back hemispheres.
    \item \textbf{Fix:} Lower \textbf{Sigmoid K} to 10 or 12 in the Simulation Settings card. This creates a wider, smoother blend region typical of physical antennas.
\end{itemize}

\subsection{Why do I see extra large lobes in my Phased Array pattern?}
\begin{itemize}
    \item \textbf{Cause:} \textbf{Grating Lobes}. This occurs when the element spacing is greater than $0.5\lambda$ (specifically when $d > \lambda / (1 + |\sin\theta_{steer}|)$).
    \item \textbf{Fix:} Reduce \textbf{Spacing X/Y} to 0.5, or reduce the steering angle.
\end{itemize}

\subsection{Why is the gain lower than Max Gain at the band edges?}
\begin{itemize}
    \item \textbf{Cause:} \textbf{VSWR Rolloff} is enabled. This feature simulates impedance mismatch loss at $f_{min}$ and $f_{max}$.
    \item \textbf{Fix:} Disable \textbf{VSWR Rolloff} in the Physics Features card to maintain constant peak gain across the entire band.
\end{itemize}

\subsection{Why does the pattern look "fuzzy" or jagged?}
\begin{itemize}
    \item \textbf{Cause:} \textbf{Noise} is set greater than 0.0, or \textbf{Pattern Breakup} is enabled.
    \item \textbf{Fix:} Set \textbf{Noise} to 0.0 and toggle \textbf{Pattern Breakup} OFF.
\end{itemize}

\newpage
\section{Parameter Fitting (Reverse Engineering)}

The Analysis page includes a \textbf{Parameter Fit} card that can derive AntennaForge configuration parameters from pre-existing measured or legacy patterns.

\subsection{How It Works}
\begin{enumerate}
    \item \textbf{Measurement Extraction:} The fitter analyses reference pattern files to extract peak gain, boresight gain, beamwidths, front-to-back ratio, and sidelobe levels.
    \item \textbf{Multi-Strategy Search:} Each run tries a different combination of physics features (sidelobes, ground reflection, etc.) with systematically varied initial parameters.
    \item \textbf{Coarse-to-Fine Optimisation:} Stage 1 evaluates on a 15\textdegree{} grid for fast exploration; Stage 2 refines on a 3\textdegree{} grid. High-dimensional types (Dish, Array) use differential evolution.
    \item \textbf{Warm-Start Cache:} Successful fits are cached at \texttt{\~{}/.antennaforge/fit\_cache.json} and used to seed future runs.
\end{enumerate}

\subsection{Per-Type Feasibility}
Not all antenna types are equally easy to fit. The table below summarises typical feasibility:

\begin{longtable}{p{3cm} p{2.5cm} p{3cm} p{5cm}}
\toprule
\textbf{Type} & \textbf{Feasibility} & \textbf{Target RMS} & \textbf{Notes} \\
\midrule
Monopole & High & 0.5\,dB & Few parameters, closed-form model. \\
Omni & High & 1.0\,dB & Elevation-only shaping; azimuth is uniform. \\
LPDA & Medium--High & 1.5\,dB & Gaussian beam is flexible; frequency scaling is the main challenge. \\
Horn & Medium--High & 1.5\,dB & Asymmetric E/H planes well-modelled. \\
Panel & Medium & 1.5\,dB & Good for sector antennas; sinc sidelobes may not match all real panels. \\
Dish & Medium--Low & 2.0\,dB & Many parameters (feed taper, aperture); needs 8+ runs. \\
Array & Low--Medium & 2.5\,dB & Highest dimensionality (element count, spacing, taper, steering); needs 10+ runs. \\
\bottomrule
\end{longtable}

\subsection{Tips for Successful Fitting}
\begin{itemize}
    \item \textbf{Use enough runs.} The recommended number is shown next to the Runs field. Set to 0 for unlimited runs and use the Stop button when satisfied.
    \item \textbf{Provide clean data.} The fitter works best with copolarisation-only patterns. Remove cross-pol files before fitting.
    \item \textbf{Start with a simpler type.} If unsure, try LPDA first --- its Gaussian model is the most forgiving.
    \item \textbf{Check the RMS.} The report shows per-file RMS error. Values below the per-type target are generally ``close enough'' for system-level work.
    \item \textbf{Use Pattern Set Comparison.} After fitting, generate patterns from the fitted config and compare them against the originals using the Pattern Set Comparison card to see per-frequency percentage differences.
\end{itemize}

\subsection{Pattern Set Comparison}
The \textbf{Pattern Set Comparison} card compares two sets of pattern files frequency-by-frequency. For each metric (peak gain, boresight gain, azimuth beamwidth, elevation beamwidth, minimum gain), it computes:
\[ \Delta\% = \frac{B - A}{|A|} \times 100 \]
This is useful for validating a fitted pattern against the original measured data, or for comparing two different antenna configurations across their operating band. The result includes per-file differences and summary statistics (average and maximum absolute difference).

\subsection{Batch Fitting}
Batch Fitting allows you to fit many pattern sets at once from a parent directory of subfolders. Instead of running the fitter manually on each antenna configuration, you point it at a top-level directory and it processes every subfolder automatically.

\begin{itemize}
    \item \textbf{Directory Structure:} Each subfolder within the parent directory represents one pattern set. The subfolder name is used as the set name in the final report.
    \item \textbf{Auto-Detection:} The batch fitter reads \texttt{generation\_config.json} (if present in the subfolder) to determine \texttt{antenna\_type}, \texttt{f\_min\_mhz}, and \texttt{f\_max\_mhz}. If no config file is found, it falls back to parsing filenames using the \texttt{TYPE\_FREQpDECMHz\_copol.csv} naming convention.
    \item \textbf{Pipeline:} For each subfolder the process is: auto-detect parameters $\rightarrow$ fit (2$\times$ recommended runs by default) $\rightarrow$ save \texttt{fitted\_config.json} $\rightarrow$ generate patterns from fitted config $\rightarrow$ compare original vs fitted $\rightarrow$ produce detailed report.
\end{itemize}

\subsubsection*{Batch Report Contents}
The output file \texttt{batch\_fit\_report.txt} contains:
\begin{itemize}
    \item \textbf{Grand Summary Table:} An overview of all sets showing antenna type, frequency band, number of files, RMS error, percentage error, comparison RMS, and processing time.
    \item \textbf{Per-Set Detailed Sections:} For each pattern set the report includes:
    \begin{itemize}
        \item Fit quality metrics (mid-band and aggregate).
        \item Per-file breakdown of fitting error.
        \item Original vs generated comparison with gain, beamwidth, and minimum-gain tables.
        \item The full fit report text for that set.
    \end{itemize}
\end{itemize}

\subsubsection*{Tips for Batch Fitting}
\begin{itemize}
    \item Organise patterns so that each subfolder contains files from one antenna or configuration.
    \item Include a \texttt{generation\_config.json} in each subfolder for reliable auto-detection of antenna type and frequency range.
    \item Use \textbf{Runs\,=\,0} for automatic run count (2$\times$ recommended), or set a higher value for difficult types like Array.
    \item Review \texttt{batch\_fit\_report.txt} to identify which sets fitted well and which need manual attention.
\end{itemize}

\newpage
\section{Glossary}

\begin{description}
    \item[dBi] Decibels relative to an isotropic radiator. $0 \text{ dBi} = 10 \log_{10}(1)$.
    \item[dBd] Decibels relative to a half-wave dipole. $0 \text{ dBd} \approx 2.15 \text{ dBi}$.
    \item[Beamwidth (HPBW)] The angular width of the main beam at the -3 dB (half-power) points.
    \item[Front-to-Back Ratio (F/B)] The difference in gain between the peak of the main beam ($0^\circ$) and the rear of the antenna ($180^\circ$).
    \item[Sidelobe] Any radiation lobe other than the main beam.
    \item[Grating Lobe] A secondary main beam caused by spatial aliasing in phased arrays with large element spacing.
    \item[Squint] A shift in the main beam direction away from the mechanical boresight, often frequency-dependent.
    \item[Cross-Polarization] Radiation orthogonal to the intended polarization (e.g., horizontal component of a vertical antenna).
    \item[VSWR] Voltage Standing Wave Ratio. A measure of impedance match. High VSWR implies reflected power and reduced radiated gain.
\end{description}

\end{document}